% !PRINTED 643
\weberpage{643}
Setzt man darin
\[
\Delta = \pm\,q^{\,q^{\varkappa-1}[\varkappa(q-1)-1]}
\]
so folgt
\weq{7}{
\begin{array}{rl}
\Delta_1 = q^{\varkappa\, q^{\varkappa-1}\frac{q-1}{2} - \frac{q^{\varkappa-1}+1}{2}} & \text{bei ungeradem } q\\
= 2^{(\varkappa-1)2^{\varkappa-2}-1} & \text{f\"ur } q=2,
\end{array}
}
so dass $\Delta_1$ immer eine Potenz von $q$ ist.

\begin{center}
\S.\,177.\\[4pt]
\gesperrt{Die Primideale im K\"orper $H_m$.}
\end{center}

Es ist jetzt die Frage zu untersuchen, in welcher Beziehung die Primideale des K\"orpers $H_m$ zu denen des K\"orpers $\Omega_m$ stehen.

Wenn $\mathfrak{P}$ irgend ein Primideal in $H_m$ vom Grade $f_1$ bedeutet, so muss $\mathfrak{P}$ Theiler einer nat\"urlichen Primzahl $p$ sein, und $\mathfrak{P}$ kann also (nach \S.\,171) nur dann durch die zweite oder eine h\"ohere Potenz eines Primfactors $\mathfrak{p}$ in $\Omega_m$ theilbar sein, wenn $p=q$ ist.

Ist aber $\mathfrak{P}$ ein Theiler von $q$, so muss es eine Potenz von $\sigma=1-r$ sein. Wir setzen etwa
\[
\mathfrak{P}=\sigma^\lambda.
\]
Nehmen wir hiervon die Norm in Bezug auf $\Omega_m$, die das Quadrat der Norm in Bezug auf $H_m$ ist, so folgt (nach \S.\,169)
\[
q^{2f_1}=q^\lambda,
\]
also $f_1=\frac12\lambda$. Da $f_1$ eine ganze Zahl ist, so muss $\lambda$ mindestens $=2$ sein. Es kann aber $\lambda$ auch nicht gr\"osser als $2$ sein, da $\sigma^2$ mit der in $H_m$ enthaltenen ganzen Zahl $(1-r)(1-r^{-1})$ associirt ist, also $\sigma^2$ durch $\mathfrak{P}$ theilbar sein muss. Demnach ist $f_1=1$, und wir erhalten den Satz:

1. \gesperrt{Die Primzahl $q$ ist die $\frac12\mu^{\text{te}}$ Potenz einer in $H_m$ existirenden Primzahl ersten Grades.}

Es sei nun $p$ von $q$ verschieden, und $\mathfrak{p}$ ein Primfactor von $\mathfrak{P}$ im K\"orper $\Omega_m$ vom Grade $f$, der durch die Substitution $(r,r^{-1})$ in $\mathfrak{p}'$ \"ubergeht. Dann ist $\mathfrak{P}$ sowohl durch $\mathfrak{p}$ als durch $\mathfrak{p}'$ theilbar, und andererseits ist $\mathfrak{p}\mathfrak{p}'$ in $H_m$ enthalten und also durch $\mathfrak{P}$ theilbar, und wir haben zu unterscheiden, ob $\mathfrak{p}, \mathfrak{p}'$ verschieden sind oder nicht.

Ist $\mathfrak{p}$ von $\mathfrak{p}'$ verschieden, so ist $\mathfrak{P}$ durch $\mathfrak{p}\mathfrak{p}'$ theilbar, und es folgt $\mathfrak{P}=\mathfrak{p}\mathfrak{p}'$, und durch Normbildung $f_1=f$.
